\documentclass[11pt]{article}
\usepackage[T1]{fontenc}
\usepackage[a4paper,margin=1in]{geometry}
\usepackage{amsmath,amssymb,amsthm,mathtools}
\usepackage{graphicx}
\usepackage[hidelinks]{hyperref}
\usepackage{microtype}
\usepackage{enumitem}

\newtheorem{theorem}{Theorem}
\newtheorem{lemma}[theorem]{Lemma}
\theoremstyle{definition}
\newtheorem*{sourcequestion}{Source Question 7.3}

\newcommand{\MIN}{\operatorname{MIN}}
\newcommand{\MAX}{\operatorname{MAX}}
\newcommand{\sa}{\mathrm{sa}}

\setlength{\parindent}{0pt}
\setlength{\parskip}{0.55em}

\title{A Positively Generated Function System with No\\
Uniform Complete Positive-Generation Bound}
\author{Run \texttt{fa\_banach\_001} --- lane 2}
\date{11 August 2026}

\begin{document}
\maketitle

\begin{center}
\fbox{\parbox{0.91\textwidth}{
\textbf{Status: candidate counterexample, likely valid, human review required.}
Question~7.3 of Humeniuk--Kennedy--Manor asks whether positive generation of a
function system forces complete $\alpha$-generation for one uniform
$\alpha$.  The explicit Cantor-cube function system below is positively
generated but has no such uniform constant, giving a negative answer.
}}
\end{center}

\section{Source question}

Humeniuk, Kennedy, and Manor study dualizability and positive decomposition
for possibly nonunital operator systems in arXiv:2312.04791, later published
in the \emph{Journal of Functional Analysis} \textbf{289} (2025), 111153
\cite{source}.  Their published Question~7.3, on arXiv PDF page~24, is:

\begin{center}
\includegraphics[width=0.97\textwidth]{figures/open_problem_crop.png}
\end{center}

\begin{sourcequestion}
If $S\subseteq C(X)$ is a positively generated function system, is $S$
completely $\alpha$-generated for some $\alpha>0$?
\end{sourcequestion}

Here a function system may be nonunital.  An ordered $*$-Banach space $E$ is
$\alpha$-generated if every $x\in E^\sa$ has a decomposition $x=y-z$ with
$y,z\in E^+$ and
\[
  \|y\|+\|z\|\leq \alpha\|x\|.
\]
A matrix ordered operator space is completely $\alpha$-generated if the same
constant works in every $M_n(E)$.

\section{Proof Intuition}

The source already constructs a matrix ordered version of $\MIN(\ell^1)$
whose positive cones are coordinatewise and whose complete positive-generation
constants are unbounded.  Its unresolved issue is that the example is not
known there to be a function system.

The key observation is that an equivalent norm makes precisely this ordered
space concrete and commutative.  On the Cantor cube $K=\{0,1\}^{\mathbb N}$,
send $z\in\ell^1$ to the continuous subset-sum function
$f_z(\varepsilon)=\sum_k\varepsilon_kz_k$.  Singleton points recover each
coefficient, so positivity at every matrix level is exactly coordinatewise
positivity.  The resulting norm is equivalent to $\|z\|_1$, and both the
concrete function system and $\MIN(\ell^1)$ are minimal operator spaces.
Thus their coefficient identification is a completely bounded complete order
isomorphism.  A uniform generation constant would transfer through it,
contradicting the source's example.

\section{Counterexample}

Let
\[
  K=\{0,1\}^{\mathbb N}
\]
with the product topology.  For $z=(z_k)\in\ell^1$ define
\[
  f_z(\varepsilon)=\sum_{k=1}^{\infty}\varepsilon_k z_k,
  \qquad \varepsilon=(\varepsilon_k)\in K,
\]
and put
\[
  S=\{f_z:z\in\ell^1\}\subseteq C(K).
\]

\begin{theorem}
The space $S$ is a closed, selfadjoint, positively generated function system.
It is not completely $\alpha$-generated for any $\alpha>0$.  Consequently,
Question~7.3 has a negative answer.
\end{theorem}

We separate the one external obstruction used in the proof.

\begin{lemma}[Humeniuk--Kennedy--Manor]
Let $X=\MIN(\ell^1)$, equipped at every matrix level with the coordinatewise
positive cone
\[
 M_n(X)^+=\{(A_k)_{k\geq1}:A_k\in M_n^+\}.
\]
Then $X$ is not completely $\beta$-generated for any finite $\beta$.
\end{lemma}

\begin{proof}
This is Example~8.6 on arXiv PDF pages~28--29 of the source \cite{source}.
For completeness, its obstruction is recalled here.  The dual matrix ordered
operator space is $X^*=\MAX(\ell^\infty)$ with the usual coordinatewise
ordering.  If $X$ were completely $\beta$-generated, the source's
Proposition~5.5 would make $X^*$ completely $2\beta$-normal.

The minimal and maximal operator-space structures on the infinite-dimensional
space $\ell^\infty$ are not completely equivalent
\cite[Theorem~14.3]{paulsen}.  Hence, for arbitrarily large $R$, there are
$n$ and $u\in M_n(\ell^\infty)$ with
\[
 \|u\|_{\min}\leq1,
 \qquad \|u\|_{\max}>4R.
\]
In the C$^*$-algebra $M_n(\ell^\infty)$, write $u$ as the linear combination
of the four positive and negative parts of $\operatorname{Re}u$ and
$\operatorname{Im}u$.  Each of these four positive elements has minimal norm
at most one, and one of them, say $y$, has maximal norm greater than $R$.
Thus
\[
 0\leq y\leq 1,
 \qquad \|1\|_{\max}=1,
 \qquad \|y\|_{\max}>R.
\]
This violates every uniform complete normality constant.  Taking
$R>2\beta$ gives the desired contradiction.
\end{proof}

\begin{proof}[Proof of the theorem]
\textbf{Step 1: $S$ is a closed function system.}
The series defining $f_z$ converges uniformly, because its tail is bounded by
the corresponding $\ell^1$ tail.  Thus $f_z\in C(K)$.  The map
$z\mapsto f_z$ is complex linear, injective (evaluate at the singleton
points), and respects conjugation.

We have the upper estimate
\[
 \|f_z\|_\infty\leq\|z\|_1.
\]
For the reverse estimate, given $\theta\in[0,2\pi]$, let
\[
 E_\theta=\{k:\operatorname{Re}(e^{-i\theta}z_k)\geq0\}
\]
and let $1_{E_\theta}\in K$ be its characteristic sequence.  Then
\[
 \bigl|f_z(1_{E_\theta})\bigr|
 \geq \sum_k\max\{\operatorname{Re}(e^{-i\theta}z_k),0\}.
\]
The average over $\theta$ of the $k$th summand is $|z_k|/\pi$.
Tonelli's theorem therefore gives some $\theta$ for which
\[
  \|f_z\|_\infty\geq\frac1\pi\|z\|_1.
\]
The coefficient map is consequently a Banach-space isomorphism, so its image
$S$ is closed.  Hence $S$ is a concrete (possibly nonunital) function system.

\textbf{Step 2: determine all matrix cones.}
If $f_z=f_z^*$, singleton evaluations show that every $z_k$ is real.
Furthermore,
\[
 f_z\geq0 \text{ on }K
 \quad\Longleftrightarrow\quad z_k\geq0\text{ for every }k.
\]
The forward implication again uses singleton points, and the reverse
implication is immediate.  Thus $f_z=f_{z^+}-f_{z^-}$ for selfadjoint $f_z$,
and $S$ is positively generated.

More generally, an element $F\in M_n(S)$ has a unique absolutely summable
coefficient sequence $(A_k)$ in $M_n$ and
\[
 F(\varepsilon)=\sum_k\varepsilon_kA_k.
\]
If $F\geq0$, evaluation at the singleton point $e_k\in K$ gives $A_k\geq0$.
Conversely, if every $A_k\geq0$, every subset sum above is positive.  Hence
\[
  M_n(S)^+\quad\text{is exactly the coordinatewise cone for every }n.
\]

\textbf{Step 3: compare completely with the obstruction space.}
Let $X=\MIN(\ell^1)$ with the coordinatewise matrix cones from the lemma, and
define the coefficient map
\[
 T:S\longrightarrow X,
 \qquad T(f_z)=z.
\]
Step~2 says that $T$ and $T^{-1}$ preserve positivity at every matrix level,
so $T$ is a complete order isomorphism.

Every concrete subspace of a commutative C$^*$-algebra has its minimal
operator-space structure \cite[Proposition~3.3.1]{paulsen}; hence $S$ is
minimal, while $X$ is minimal by definition.  The universal property of
minimal operator spaces says that every bounded map into a minimal operator
space is completely bounded with the same norm \cite{blecher}.  The scalar
estimates from Step~1 therefore imply
\[
 \|T\|_{\mathrm{cb}}\leq\pi,
 \qquad \|T^{-1}\|_{\mathrm{cb}}\leq1.
\]

\textbf{Step 4: transfer a hypothetical uniform constant.}
Suppose that $S$ were completely $\alpha$-generated.  Given
$x\in M_n(X)^\sa$, put $s=T_n^{-1}x$ and write $s=p-q$ with $p,q\geq0$ and
\[
 \|p\|+\|q\|\leq\alpha\|s\|.
\]
Then $x=T_np-T_nq$ is a positive decomposition in $M_n(X)$ and
\[
 \|T_np\|+\|T_nq\|
 \leq \|T\|_{\mathrm{cb}}\alpha\|T^{-1}\|_{\mathrm{cb}}\|x\|
 \leq \pi\alpha\|x\|.
\]
The same constant works for every $n$, so $X$ would be completely
$\pi\alpha$-generated.  This contradicts the lemma.  Therefore no finite
complete generation constant exists for $S$.
\end{proof}

\section{Verification}

The reusable script \texttt{code/verifier.py} exhaustively evaluates all
subsets in seeded finite-coordinate examples.  It checks the exact real norm
formula, the complex $1/\pi$ lower bound, and the matrix singleton-cone
mechanism.  Run it with
\begin{verbatim}
conda run --no-capture-output -n sandbox python \
  runs/fa_banach_001/solutions/counterexamples/\
  2312.04791_cantor_cube_positive_generation_no_uniform_complete_bound/\
  code/verifier.py
\end{verbatim}
These finite checks are sanity tests, not a proof of the infinite-dimensional
MIN/MAX obstruction.  The complete proof is the argument above.  The separate
\texttt{VERIFIER\_REPORT.md} records the review checklist and output.

\section{Novelty, limitations, and review recommendation}

A bounded search through 11 August 2026 checked the run's four lightweight
indexes, the exact wording of published Question~7.3, the phrases
``positively generated function system'', ``completely alpha-generated'',
and ``MIN($\ell^1$) function system'', and later sources citing
arXiv:2312.04791.  The 2025 JFA version still states the question as open.
No later answer or matching Cantor-cube construction was found.  Novelty is
therefore plausible but not certified.

The counterexample is nonunital, exactly within the source's convention that
function systems may be nonunital.  It does not address a strengthened unital
variant.  It also does not quantify the optimal generation constant at each
fixed matrix level.

Human review is recommended, with priority on two standard but decisive
points: the complete-order identification of the coordinate cones, and the
minimal-operator-space universal property used to promote scalar norm
equivalence to cb norm equivalence.

\begin{thebibliography}{9}
\bibitem{source}
A. Humeniuk, M. Kennedy, and N. Manor,
\emph{An extension property for noncommutative convex sets and duality for
operator systems}, J. Funct. Anal. \textbf{289} (2025), 111153;
arXiv:2312.04791.

\bibitem{paulsen}
V. Paulsen,
\emph{Completely Bounded Maps and Operator Algebras},
Cambridge Studies in Advanced Mathematics 78, Cambridge University Press,
2002.

\bibitem{blecher}
D. P. Blecher,
\emph{Introduction to Operator Spaces}, lecture notes, Fields Institute;
minimal-operator-space universal property in Section~1.2.
\end{thebibliography}

\end{document}
